\section{Inyección de Dependencias}

\subsection{Regla: inyección por constructor y dependencia de abstracciones}

\textbf{Regla:} Toda clase que dependa de un servicio externo (por ejemplo, un generador de números aleatorios, un registro de partidas o un servicio de persistencia) deberá recibir esa dependencia a través de su constructor, declarada mediante una interfaz, en lugar de instanciar la implementación concreta dentro de la propia clase.

\textbf{Descripción:} Esta regla es deliberadamente independiente de cualquier contenedor de inyección de dependencias concreto y de cualquier framework: aplica igual sin importar qué mecanismo de resolución de dependencias decida adoptar el equipo, o si no se adopta ninguno y las dependencias se ensamblan manualmente. Martin formula este principio como la inversión de dependencias: los módulos de alto nivel no deben depender de módulos de bajo nivel, sino que ambos deben depender de abstracciones (Martin, 2008).

\textbf{Código correcto:}
\begin{lstlisting}[style=csharpstyle]
public sealed class TurnManager
{
    private readonly IDiceRoller _diceRoller;

    public TurnManager(IDiceRoller diceRoller)
    {
        _diceRoller = diceRoller;
    }

    public int PlayTurn()
    {
        return _diceRoller.Roll();
    }
}
\end{lstlisting}

\textbf{Código incorrecto:}
\begin{lstlisting}[style=csharpstyle]
public sealed class TurnManager
{
    private readonly DiceRoller _diceRoller;

    public TurnManager()
    {
        _diceRoller = new DiceRoller();
    }

    public int PlayTurn()
    {
        return _diceRoller.Roll();
    }
}
\end{lstlisting}
